% ICLR appendix fragment: synthetic observation reconstruction and robustness.
% Required packages in the parent document: amsmath, graphicx.
% Paths assume compilation from the repository root and may be redefined.
% The two single-dataset audit PDFs are copied into the robustness directory
% so that this appendix section has one self-contained figure bundle.
\providecommand{\syntheticindividualauditfigdir}{analyses/artifacts/synthetic/prediction_robustness}
\providecommand{\syntheticpredictionrobustnessfigdir}{analyses/artifacts/synthetic/prediction_robustness}

\subsection{Observation reconstruction under single- and multi-dataset training}
\label{app:synthetic_prediction_robustness}

This section asks a deliberately narrower question than recovery of the
programmed kinetic profile $K_t$: does joint training on increasingly many
biased observations preserve the model's ability to reconstruct each
dataset-specific observed profile?  This distinction is essential.  The
fitted mean factorizes as
\begin{equation}
    \mu_{dti}=S_{dt}L_{ti}\gamma_{dti},
    \label{eq:appendix_mu_factorization}
\end{equation}
so an accurate $\mu_{dt}$ can coexist with an incorrect allocation of
structure between the shared profile $L_t$ and the dataset correction
$\gamma_{dt}$.  Observation reconstruction is therefore a capacity and
robustness diagnostic, not evidence by itself for successful disentanglement.

\paragraph{Notation and evaluation target.}
Let $d$ index one of the ten programmed bias conditions, $t$ a held-out
transcript, $r\in\{1,2\}$ a sampled replica, and $i$ a model-valid exported
profile position.  The observed target supplied to the model is the arithmetic
two-replica consensus
\begin{equation}
    \overline{Y}_{dti}=\frac{1}{2}
       \left(Y^{(1)}_{dti}+Y^{(2)}_{dti}\right).
    \label{eq:appendix_replica_consensus}
\end{equation}
For a cumulative model trained on the ordered panel $\mathcal D_N$ of $N$
conditions, we calculate a Pearson correlation separately within every
transcript--dataset profile,
\begin{equation}
    r^{(N)}_{dt}
      =\operatorname{PCC}_{i\in\mathcal M_{dt}}
         \!\left(\mu^{(N)}_{dti},\overline{Y}_{dti}\right),
    \qquad d\in\mathcal D_N,
    \label{eq:appendix_mu_observation_pcc}
\end{equation}
where $\mathcal M_{dt}$ is the finite exported model mask.  Correlations are
not pooled across codons or transcripts, and no Fisher transform, smoothing,
zero filtering, or reliability weighting is applied during evaluation.
Undefined correlations from constant profiles remain missing rather than
being set to zero.  To match the raw validation diagnostic and the existing
single-dataset audit exactly, the exported mask retains the validated appended
zero terminal position; no additional boundary crop is introduced here.

All comparisons below use training seed 42, uniform reference weights
$\pi_d=1/N$, and the frozen checkpoint minimizing validation loss.  Selecting
this checkpoint rather than the checkpoint maximizing validation PCC avoids
directly selecting on the displayed statistic.  The same 1,920 validation
transcripts are used by every bias and every $N$ within a read depth.  No
profile is undefined in the complete grid, so the fixed complete cohort
contains all 1,920 transcripts at each depth.  The cohorts differ across
depths and have SHA-256 prefixes \texttt{7f02618d4975},
\texttt{361ae47698c2}, and \texttt{1a508cab4970} at 0.25, 2, and 20 reads per
codon, respectively.  Consequently, the three depths are presented as
parallel sensitivity analyses rather than silently treated as if they used
identical transcripts.

\paragraph{Single-dataset prediction and empirical repeatability.}
For the independently trained single-dataset models, we compare
$r^{(1)}_{dt}$ with the within-dataset replica correlation
\begin{equation}
    a_{dt}=\operatorname{PCC}_{i\in\mathcal M_{dt}}
       \!\left(Y^{(1)}_{dti},Y^{(2)}_{dti}\right).
    \label{eq:appendix_replica_pcc}
\end{equation}
Figure~\ref{fig:synthetic_single_prediction_replica_boxes} shows the full
matched transcript distributions.  Across the ten biases, the median of the
condition-level median replica PCCs increases from 0.171 to 0.535 and 0.736
with read depth.  The corresponding medians for
$\operatorname{PCC}(\mu,\overline Y)$ are 0.538, 0.831, and 0.917.  The latter
quantity may exceed replica--replica agreement because the target averages
the same two replicas and the sequence model shares statistical strength
across training transcripts.  Replica agreement is therefore not a formal
upper bound or an independent noise ceiling for model--consensus agreement.

\begin{figure*}[t]
    \centering
    \includegraphics[width=0.94\textwidth]
      {\syntheticindividualauditfigdir/single_dataset_prediction_vs_replica_agreement_boxplots.pdf}
    \caption{\textbf{Single-dataset prediction versus replica agreement.}
    Gray boxes show the transcript-level distribution of
    $\operatorname{PCC}(Y^{(1)}_{dt},Y^{(2)}_{dt})$; purple boxes show
    $\operatorname{PCC}(\mu^{(1)}_{dt},\overline Y_{dt})$ for the matching
    depth, bias, and held-out transcript.  Boxes span the interquartile range,
    central lines are medians, and whiskers span the 5th--95th percentiles.
    Each condition contains 1,920 transcripts.  These are dependent
    comparisons because $\overline Y$ is formed from the same two replicas.}
    \label{fig:synthetic_single_prediction_replica_boxes}
\end{figure*}

The condition ordering is nevertheless highly consistent.  Define
\begin{equation}
    \rho_C
      =\operatorname{Spearman}_{d}
       \left(
          \operatorname{median}_{t}a_{dt},
          \operatorname{median}_{t}r^{(1)}_{dt}
       \right).
    \label{eq:appendix_condition_spearman}
\end{equation}
Across the ten bias conditions, $\rho_C$ is 0.976, 1.000, and 0.952 at 0.25,
2, and 20 reads per codon (Figure~\ref{fig:synthetic_single_prediction_spearman}).
Thus conditions whose two sampled replicas are easier to reconcile also tend
to be those whose consensus is easier for the single-dataset model to predict.
This is a descriptive association across ten engineered conditions, not a
transcript-level correlation or an inferential estimate for experimental
Ribo-seq datasets.

\begin{figure*}[t]
    \centering
    \includegraphics[width=\textwidth]
      {\syntheticindividualauditfigdir/single_dataset_prediction_vs_replica_agreement_scatter.pdf}
    \caption{\textbf{Condition-level association between repeatability and
    single-dataset prediction.}  Each point is one programmed bias condition.
    The horizontal coordinate is its median transcript replica PCC and the
    vertical coordinate is its median transcript
    $\operatorname{PCC}(\mu,\overline Y)$.  The dashed line denotes equality;
    the reported Spearman coefficient ranks the ten condition medians.  It
    does not test independence between the two axes and should not be read as
    a noise-ceiling analysis.}
    \label{fig:synthetic_single_prediction_spearman}
\end{figure*}

\paragraph{Matched single-versus-cumulative comparison.}
The cumulative panels use the fixed order $3'$-AA, $3'$-CC, $3'$-GG,
$3'$-UU, $5'$-AA, $5'$-CC, $5'$-GG, $5'$-UU, GC-rich, and AU-rich.  A raw
panel average at $N$ is confounded by this changing composition.  We therefore
pair every cumulative score with the independently fitted single-dataset model
for the same depth, named bias, and transcript,
\begin{equation}
    \Delta^{(N)}_{dt}=r^{(N)}_{dt}-r^{(1,d)}_{dt},
    \qquad d\in\mathcal D_N,
    \label{eq:appendix_paired_mu_delta}
\end{equation}
and summarize both models on the same prefix $\mathcal D_N$:
\begin{align}
    R^{(N)}
      &=\frac{1}{|\mathcal T_C|}\sum_{t\in\mathcal T_C}
        \frac{1}{N}\sum_{d\in\mathcal D_N}r^{(N)}_{dt},\\
    R^{(1\mid N)}
      &=\frac{1}{|\mathcal T_C|}\sum_{t\in\mathcal T_C}
        \frac{1}{N}\sum_{d\in\mathcal D_N}r^{(1,d)}_{dt}.
    \label{eq:appendix_panel_reconstruction_means}
\end{align}
Here $R^{(1\mid N)}$ changes slightly with $N$ because it is the matched mean
of the first $N$ bias-specific fits, not one privileged $N=1$ model.  This
construction prevents an easy or difficult newly added bias from being
misreported as a benefit or failure of joint training.

Figure~\ref{fig:synthetic_mu_robustness_aggregate} shows that joint training
preserves observation reconstruction.  At the two endpoints $N=2$ and
$N=10$, respectively, cumulative mean PCC is 0.541 and 0.547 at 0.25
reads/codon, 0.819 and 0.822 at 2 reads/codon, and 0.908 and 0.909 at 20
reads/codon.  At $N=10$, the paired cumulative-minus-single
differences are $+0.00162$ (95\% transcript-bootstrap interval
$[+0.00140,+0.00183]$), $+0.00041$ ($[+0.00028,+0.00054]$), and $-0.00014$
($[-0.00026,-0.00002]$), respectively.  Their small magnitude is more
important than their sign: there is no practically meaningful observation-fit
penalty from fitting the ten conditions jointly in this realization.

\begin{figure*}[t]
    \centering
    \includegraphics[width=\textwidth]
      {\syntheticpredictionrobustnessfigdir/synthetic_mu_robustness_aggregate.pdf}
    \caption{\textbf{Observed-profile reconstruction is stable under
    cumulative multi-dataset training.}  \textbf{(A)} Solid curves show
    $R^{(N)}$ for the cumulative model.  Dashed curves show the matched
    $R^{(1\mid N)}$ average from the same first $N$ bias-specific
    single-dataset models.  \textbf{(B)} Mean paired difference
    $R^{(N)}-R^{(1\mid N)}$; the horizontal line denotes no change.  Error bars
    are pointwise 95\% percentile intervals from 2,000 paired transcript
    bootstrap draws.  The same 1,920-transcript cohort is retained across all
    $N$ within each depth.  Intervals condition on one set of seed-42 fitted
    models and do not represent training-seed uncertainty.}
    \label{fig:synthetic_mu_robustness_aggregate}
\end{figure*}

\paragraph{Condition-resolved audit.}
The aggregate comparison can conceal a condition-specific failure.  We
therefore retain every valid cell of the triangular design: the $d$th bias is
shown only from the first cumulative panel that contains it.  In
Figure~\ref{fig:synthetic_mu_condition_absolute}, the $N=1$ column contains the
ten independently fitted bias-specific models, whereas columns $N\geq2$
contain the corresponding rows from the cumulative model.  The absolute
scores expose the dominant depth and bias effects: cell means range from
0.443--0.641, 0.762--0.874, and 0.881--0.930 across the three depths.

\begin{figure*}[p]
    \centering
    \includegraphics[height=0.89\textheight]
      {\syntheticpredictionrobustnessfigdir/synthetic_mu_pcc_condition_resolved.pdf}
    \caption{\textbf{Condition-resolved observation reconstruction for every
    available training configuration.}  Rows are programmed bias conditions,
    columns are training configurations, and cells report the mean of 1,920
    transcript-level $\operatorname{PCC}(\mu_{dt},\overline Y_{dt})$ values.
    The first column consists of ten distinct single-dataset fits; each later
    column is one cumulative fit and therefore has entries only for its first
    $N$ included conditions.  Gray cells were not trained and are not missing
    measurements.  A common color scale is used across read depths.}
    \label{fig:synthetic_mu_condition_absolute}
\end{figure*}

Figure~\ref{fig:synthetic_mu_condition_delta} removes the large depth and bias
main effects by displaying the matched difference in
\cref{eq:appendix_paired_mu_delta}.  Across every displayed condition and
panel size, mean differences range from $+0.00041$ to $+0.00453$ at depth
0.25, from $-0.00053$ to $+0.00167$ at depth 2, and from $-0.00094$ to
$+0.00082$ at depth 20.  Hence the aggregate stability is not produced by a
large gain in some conditions cancelling a large loss in others.  The largest
cellwise changes remain below 0.005 PCC.

\begin{figure*}[p]
    \centering
    \includegraphics[height=0.89\textheight]
      {\syntheticpredictionrobustnessfigdir/synthetic_mu_pcc_delta_from_single.pdf}
    \caption{\textbf{Condition-resolved change relative to matched
    single-dataset fits.}  Each cell reports the mean paired difference
    $\Delta^{(N)}_{dt}$ after matching read depth, named bias, transcript, and
    observed target.  Positive values favor the cumulative model and negative
    values favor the separately trained single-dataset model.  The diverging
    scale is centered at zero and shared across depths.  These values quantify
    preservation of $\mu$ reconstruction; they do not compare $L_t$ with
    $K_t$ or $\gamma_{dt}$ with the injected bias.}
    \label{fig:synthetic_mu_condition_delta}
\end{figure*}

\paragraph{Interpretation and limitations.}
The two analyses support complementary conclusions.  Single-dataset
prediction quality tracks how reproducibly a condition was sampled, while
joint training preserves essentially the same condition-specific observation
fit as ten separate models.  This is the desired robustness property for the
full fitted mean: adding heterogeneous datasets does not force a trade-off in
which one observed profile is reconstructed at the expense of another.

The result should not be overstated.  First, the validation profiles are used
both for checkpoint selection through validation loss and for this diagnostic;
they are not an independent test set.  Second, there is one training seed, so
the narrow transcript-bootstrap intervals quantify transcript variation
conditional on the fitted models rather than optimization variability.
Third, increasing $N$ simultaneously increases the amount of training data
and changes the bias composition.  Finally, because only the product in
\cref{eq:appendix_mu_factorization} is evaluated, the near-zero deltas cannot
show that $L_t$ recovered the programmed kinetic target or that
$\gamma_{dt}$ recovered the injected multiplier.  Those claims require the
separate gauge-matched recovery analyses.

% Reproducible sources:
% analyses/analyze_synthetic_prediction_robustness.py
% analyses/artifacts/synthetic/prediction_robustness/condition_resolved_summary.csv
% analyses/artifacts/synthetic/prediction_robustness/panel_aggregate_summary.csv
% analyses/artifacts/synthetic/prediction_robustness/per_transcript_paired_metrics.parquet
% analyses/artifacts/synthetic/prediction_robustness/fixed_cohort_ids.csv
% analyses/artifacts/synthetic/prediction_robustness/provenance.json
